\documentclass[11pt]{article}

\usepackage{fullpage}
\usepackage{amsmath}
\usepackage{epsfig}
\usepackage{graphics}
\usepackage{caption}
\usepackage{circuitikz}

\oddsidemargin -0.15in
\evensidemargin -0.15in
\headheight 0in
\headsep 0in
%\topmargin -0.25in
\textheight 9.3in
\textwidth 6.8in

\begin{document}


\begin{center}
\section*{
Franklin \& Marshall College \hfill PHY 223\\
\bigskip
{\Large \it Speed of light}}
\end{center}

The speed of light in free space is an important and intriguing constant of
nature. Whether the light comes from a laser on a desktop or from a distant
galaxy that is hurtling away at fantastic speeds, the measurement of the speed
of light will yield the same constant value. In more precise terminology, the
speed of light is independent of the relative velocities of the light source and
the observer.

As Einstein first presented in his special theory of relativity, the speed of
light is critically important in some surprising ways:

\begin{enumerate}
 \item The speed of light establishes an upper limit to the speed that may be 
 imparted to any material object.
 
 \item Objects moving near the speed of light follow a set of physical laws 
 drastically different, not only from Newton’s Laws, but from the basic 
 assumptions of human intuition.
\end{enumerate}

It is not surprising that a great deal of time and effort has been invested in
measuring the speed of light. Some of the most accurate measurements were made
by Albert Michelson between 1926 and 1929. Michelson measured the speed of light
in air to be 2.99712x10$^8$ m/s. From this result he deduced the speed of light
in free space to be 2.99796x10$^8$ m/s. Today, the speed of light in free space
is, by definition, $(2.997924562 \pm 0.000000011)x10^8$ m/s exactly.

Our approach to the measurement of the speed of light will be direct and
intuitive, and simpler then Michelson's experiment. We will use a diode laser as
our light source. We will send pulses of light across the lab, measure how long
they take to come back, and deduce the speed of light from the knowledge that
light travels at constant speed.

\vspace{5mm}

\noindent {\bf Laser light safety notes}

We will use a relatively safe, low power laser. Nevertheless, take the following
precautions:

\begin{enumerate}

\item Never look directly into the laser beam, either directly, or as it is
reflected from a mirror or other surface.

\item Keep the optical path at floor level and keep your eyes well above floor
level at all times.

\item Only people involved with the experiment should be nearby.

\end{enumerate}

\vspace{5mm}

\noindent {\bf Experimental setup}

\begin{enumerate}
  
\item Place the breadboard with the laser and photodiode detector on the floor.
There should be at least 8-12 meters of space in front of the laser.

\item The laser diode has two connections: the power supply that connects to an
outlet and the trigger signal that is sent to the oscilloscope. The trigger
signal cable should have a 50-Ohm at the oscilloscope.

\item Connect the photodiode output to channel 2 on the oscilloscope.
  
\item Align the laser so that it propagates along a line on the floor and that
its height is constant throughout. It will make alignment easier afterwards.
Once this alignment is done,

\item Place the mirror approximately 1 m in front of the laser and use it to
reflect the beam on the light receiver sensor.

\item Check the signal on the oscilloscope. You should trigger on channel 1.
On channel 2, you should see a strong (100s of mV) that is stable in amplitude
and time.
  
\end{enumerate}

\vspace{5mm}

\noindent {\bf Procedure}

\begin{itemize}
\item Using a tape measure, record the total path length of the light pulses.

\item Using the cursors on the oscilloscope, measure the time delay between the
two traces. One of the cursors should be set to a fixed feature on the triggering
signal (channel 1). The other cursor should be set to the onset of the photodiode
signal.

\item Also using the cursors, also record your uncertainty in the time delay.

\item Move the mirror back about 60 cm and repeat the measurement.
\end{itemize}


\vspace{5mm}

\noindent {\bf Analysis}

Make a graph of the time delay as a function of round trip distance. Find the
slope of this graph and its uncertainty. Convert this slope (and uncertainty)
into a value for the speed of light (and uncertainty).


% \bibliographystyle{plain}

% \begin{thebibliography}{10}

% \bibitem {tag} L. J. van der Pauw, ``A Method of Measuring Specific
% Resistivity and Hall Effect of Discs of Arbitrary Shape,'' Philips Research
% Reports 13, 1 (1958).

% \end{thebibliography}


\end{document}

%%% Local Variables:
%%% mode: latex
%%% TeX-master: t
%%% End:
